\documentclass[10pt,a4paper]{amsart}
\usepackage[margin=19mm]{geometry}
\usepackage{amsmath,amssymb,mathtools}
\usepackage[scr=rsfs,cal=euler]{mathalfa}
\usepackage{xfrac}
\usepackage[unicode,hidelinks,bookmarks=false]{hyperref}
\newcommand{\ie}{i.e.\ }
\newcommand{\cf}{cf.\ }
\newcommand{\resp}{respectively\ }
\newcommand{\Subsection}{\subsection}
\providecommand{\nobreakdash}{\nobreak\hbox{-}}
\setlength{\emergencystretch}{2em}
\multlinegap0pt
\usepackage{bm}
\usepackage{bbm}
\usepackage{dsfont}
\usepackage{xspace}
\usepackage{cancel}
\usepackage{mathtools}
\usepackage{amsthm}
\makeatletter
\@ifundefined{theorem}{\newtheorem{theorem}{Theorem}}{}
\@ifundefined{lemma}{\newtheorem{lemma}[theorem]{Lemma}}{}
\makeatother
\title[Borel-Hirzebruch type formula for the graph equivariant coho]{Borel-Hirzebruch type formula for the graph equivariant cohomology of a projective bundle over a GKM-graph}
\author{Shintaro Kuroki, Grigory Solomadin}
\date{Source: arXiv:2207.11380v4 (CC BY 4.0)}
\begin{document}
\maketitle

\noindent\emph{Source and licence.} Borel-Hirzebruch type formula for the graph equivariant cohomology of a projective bundle over a GKM-graph. Version arXiv:2207.11380v4 (CC BY 4.0), distributed under the Creative Commons Attribution 4.0 International licence, as recorded in the arXiv licence metadata for this exact version and shown on the abstract page https://arxiv.org/abs/2207.11380v4. Full rights notice: licenses/arxiv-2207.11380-CC-BY-4.0.txt.\par\smallskip

\noindent\textit{Editorial scope.} Counted unit: the Lemma whose source label is \texttt{technical\_lemm} in the article (source0.tex lines 880--887, source label \texttt{technical\_lemm}), delivered together with the article's own complete original proof of it (source0.tex lines 888--945, one \texttt{proof} environment of 58 lines). No mathematics is written, repaired or re-derived: statement and proof are the article's own text line for line. Required context retained uncounted: none. Every definition, display and label of those lines is retained unchanged. Omitted from this package and not redistributed: 1--879, 946--1706. Nothing inside a retained excerpt is omitted or altered: every retained body is delivered exactly as the article's lines are written, so no omission receipt of any kind accompanies this package. Citations: the retained excerpts carry no citation command at all, so no bibliography entry of the article is transcribed and no reference is redistributed. Reference adaptation: none was needed and none was made; every reference inside the delivered unit resolves inside it, so no unresolved reference or citation is produced. Printed slips kept literal: no printed word, symbol, hypothesis or number of the counted statement or of its proof was altered. Layout and class: the article's own class is \texttt{amsart} with packages including babel, mathrsfs, mathtools, dsfont, url, booktabs, amsfonts, nicefrac, microtype, lipsum, enumitem, a4wide, color, amssymb; the wrapper is the lane's pinned \texttt{amsart} editorial wrapper at 10pt on A4, which restates the theorem-like environments the retained text needs and adds the article's own macro definitions that the retained text needs (none), with extra wrapper packages (bm, bbm, dsfont, xspace, cancel, mathtools). The delivered numbering is the unit's own. Assets: no retained span contains a figure environment, a graphics-inclusion command, a PDF-page inclusion, a table or a TikZ picture; no image file is copied into this package, the package declares no asset dependency and no image file is redistributed with it. Licence: version arXiv:2207.11380v4 of this article is distributed under the Creative Commons Attribution 4.0 International licence, as recorded in the arXiv licence metadata for this exact version and shown on the abstract page https://arxiv.org/abs/2207.11380v4. A full rights notice is delivered at licenses/arxiv-2207.11380-CC-BY-4.0.txt. Characters: every retained excerpt is pure ASCII, so the delivered engine, XeLaTeX with its default Latin Modern text font, reports no missing character.
\begin{lemma}
\label{technical_lemm}
As an $H^{*}(BT)$-module, $H^{*}(\Pi_{p}(\xi))$ is generated by $\{1, t_{p},t_{p}^{2},\ldots, t_{p}^{r}\}$, where  $t_{p}:=t|_{\mathcal{V}_{p}^{P(\xi)}}$ is the well-defined restriction of $t=c_{\xi}\in H^{*}(\Pi(\xi))$ to the subgraph $P_{p}(\xi)$.
Namely, for every element $X\in H^{*}(\Pi_{p}(\xi))$, there exist polynomials $Q_{0}(p),Q_{1}(p),\ldots, Q_{r}(p)\in H^{*}(BT^{n})\subset H^{*}(\Pi_{p}(\xi))$ such that
\begin{align}\label{eq:module_dec}
X=\sum_{s=0}^{r}Q_{s}(p)t_{p}^{s}.
\end{align}
\end{lemma}

\begin{proof}
Take an element $X\in H^{*}(\Pi_{p}(\xi))$.
Recall $\mathcal{V}_{p}^{P(\xi)}=[r+1]_{p}=\{(p,j)\ |\ j\in [r+1]\}$.
Consider the monomorphism $\iota\colon H^{*}(BT^{n})\to H^{*}(\Pi_{p}(\xi))$ defined by the constant functions.
One has  $X(p,1)\in {\rm Im}~\iota\subset H^{*}(\Pi_{p}(\xi))$ because of $X(p,1)\in H^{*}(BT^{n})$.

We first put $X_{1}:=X-X(p,1)$.
Then the element $X_{1}\in H^{*}(\Pi_{p}(\xi))$ satisfies $X_{1}(p,1)=0$.
By definition of the fiber $\Pi_{p}(\xi)$,
for every $(p,j)$ ($j=2,\ldots, r+1$), it follows from the congruence relations that one has
\begin{align*}
X_{1}(p,j)\equiv 0 \mod \xi_{p}^{1}-\xi_{p}^{j}=t_{p}(p,1)-t_{p}(p,j).
\end{align*}
In other words, for every $j=2,\ldots, r+1$
there exists an element $Y_{1}(p,j)\in H^{*}(BT^{n})$ such that
\begin{align*}
X_{1}(p,j)
=Y_{1}(p,j)\bigl(t_{p}(p,1)-t_{p}(p,j)\bigr).
\end{align*}
We next take the following element in $H^{*}(\Pi_{p}(\xi))$:
\begin{align*}
X_{2}:=X_{1}-Y_{1}(p,2)\bigl(t_{p}(p,1)-t_{p}\bigr),
\end{align*}
where we regard $Y_{1}(p,2), t_{p}(p,1)\in {\rm Im}~\iota\subset H^{*}(\Pi_{p}(\xi))$.
This element satisfies the equalities  $X_{2}(p,1)=X_{2}(p,2)=0$.
So, by the congruence relations, we have
\begin{align*}
X_{2}(p,j)\equiv 0 \mod t_{p}(p,l)-t_{p}(p,j),
\end{align*}
for $j=3,\ldots r+1$ and $l=1,2$.
Therefore, there exists $Y_{2}(p,j)\in H^{*}(BT^{n})$ for $j=3,\ldots, r+1$ such that
\begin{align*}
X_{2}(p,j)=Y_{2}(p,j)\prod_{l=1}^{2}\bigl(t_{p}(p,l)-t_{p}(p,j)\bigr).
\end{align*}
Note that $t_{p}(p,1)-t_{p}(p,j)$ and $t_{p}(p,2)-t_{p}(p,j)$ are linearly independent for any  $j=3,\dots,r+1$ because $\Pi(\xi)$ is a GKM graph.
Similarly, if $X_{k-1}\in H^{*}(\Pi_{p}(\xi))$ satisfies $X_{k-1}(p,l)=0$ for $l=1,\ldots, k-1$, then it follows from the congruence relations that there exists  $Y_{k-1}(p,j)\in H^{*}(BT^{n})$ for $j=k,\ldots,r+1$ such that
\begin{align*}
X_{k-1}(p,j)=Y_{k-1}(p,j)\prod_{l=1}^{k-1}\bigl(t_{p}(p,l)-t_{p}(p,j)\bigr).
\end{align*}
Therefore, if we put
\begin{align*}
X_{k}:=X_{k-1}-Y_{k-1}(p,k)\prod_{l=1}^{k-1}\bigl(t_{p}(p,l)-t_{p}\bigr)\in H^{*}(\Pi_{p}(\xi)),
\end{align*}
then one has the equality $X_{k}(p,l)=0$ for $l=1,\ldots, k$.

Put $Z_{k-1}:=Y_{k-1}(p,k)\prod_{l=1}^{k-1}\bigl(t_{p}(p,l)-t_{p}\bigr)$.
Then, $Z_{k-1}$ is an element generated by $\{1,t_{p},\ldots, t_{p}^{k-1}\}$.
Inductively, we can make
\begin{align*}
X_{r+1}=X_{r}-Z_{r}=0\in H^{*}(\Pi_{p}(\xi)).
\end{align*}
By the constructions as above, we have the equalities
\begin{align*}
X_{r+1}=X-(X(p,1)+Z_{1}+\cdots +Z_{r})=0.
\end{align*}
Since $(X(p,1)+Z_{1}+\cdots +Z_{r})$ is an element generated by $\{1,t_{p},\ldots, t_{p}^{r}\}$,
this proves the lemma.
\end{proof}

\end{document}
